% test-pdfprim's primitives in a document that loads hyperref and leaves
% \pdfoutput unset. A heading's text
% goes to the table of contents and to hyperref's bookmarks, and a
% caption's to the list of figures.
\documentclass{article}
\usepackage[colorlinks]{hyperref}
\newcommand\blue[1]{\pdfliteral{0 0 1 rg}#1\pdfliteral{0 g}}
\newcommand\turn[1]{\pdfsave\pdfsetmatrix{1 0 0 1}\rlap{#1}\pdfrestore\phantom{#1}}
\begin{document}
\tableofcontents
\listoffigures

\section{Head \blue{x}\pdfliteral direct{0 g} \turn{y}
  \pdfcolorstack0 push{0 0 1 rg}z\pdfcolorstack0 pop}\label{head}

Literal: before\pdfliteral page{0 0 1 rg} blue\pdfliteral{0 g} after,
a link to section~\ref{head} and to \href{https://example.com/}{a site}.

Image: before\pdfximage width 1cm {test-picture.png}\pdfrefximage\pdflastximage{} after.

Objects: before\pdfobj reserveobjnum\xdef\myobj{\the\pdflastobj}%
\immediate\pdfobj useobjnum\myobj{<< /X 1 >>}\pdfrefobj\myobj{} after.
Number: \ifnum\myobj>0 set\else unset\fi.

Link: before\pdfstartlink goto name{section.1}to the head\pdfendlink{} after.

Stack: before\pdfcolorstack0 push{0 0 1 rg}pushed\pdfcolorstack0 pop{} after.

Footnote: before\footnote{Note \blue{six} \turn{end}.} after.

\begin{figure}[h]
\centering Figure text.
\caption{Caption \blue{seven} \turn{end}}
\end{figure}

\end{document}
